\section[Campaña geométrico--topológica]
{Campaña dinámica geométrico--topológica: protocolo y niveles de evidencia}
\label{sec:campana-gt}

La campaña combina localizadores armónicos y geométricos con integración
causal, análisis de cuencas y validación topológica. Sus criterios se fijan
antes de clasificar las trayectorias y responden a cuatro preguntas:

\begin{enumerate}
  \item ¿la semilla produce una evolución causal acotada durante la ventana
        observada?;
  \item ¿esa evolución converge hacia un destino ya clasificado o revela un
        conjunto recurrente distinto?;
  \item ¿el destino y la frontera de su cuenca sobreviven a refinamientos del
        integrador, del horizonte, de la memoria y del banco de semillas?;
  \item ¿existe una certificación topológica que permita pasar de evidencia
        finita a una afirmación matemática dentro de un dominio declarado?
\end{enumerate}

\begin{hallazgo}{Estado de la campaña B0--B2}
La campaña B0--B2 se ejecutó para PLL, MAVPD y Wu, con 24 semillas,
90 trayectorias causales y 12 evaluaciones de borde. El máximo alcanzado es
EV--TG3 para una ruta del PLL y EV--TG2 para MAVPD y Wu; los detalles aparecen
en la \cref{sec:resultados-piloto-gt}. Fronteras de cuenca certificadas,
trayectorias prolongadas de borde, bloques aislantes, índices de Conley y
decisiones de ocultedad permanecen \textbf{pendientes}.
\end{hallazgo}


\subsection{Definición común del sistema y del espacio evaluado}
\label{camp:contrato-comun}

Cada experimento fijará el campo \(f\), el
vector de parámetros, el dominio físico, las unidades o la
adimensionalización, la regularidad por regiones y el grupo de simetrías. En
el caso ordinario se estudiará
\begin{equation}
 \dot x=f(x;\mu),\qquad x(0)=x_0,
 \label{camp:eq-ode-contract}
\end{equation}
mientras que la ruta fraccionaria conmensurable fijará explícitamente
\begin{equation}
 {}^{\mathrm C}D_{t_0+}^{q}x(t)=f(x(t);\mu),
 \qquad 0<q<1,\qquad x(t_0)=x_0.
 \label{camp:eq-caputo-contract}
\end{equation}
Si un caso usa órdenes no conmensurables, se especificará cada \(q_i\) y
conservará el vector de órdenes durante la evaluación.

Para comparar distancias entre variables con escalas diferentes se fijará una
matriz diagonal \(S=\operatorname{diag}(s_1,\ldots,s_n)\), \(s_i>0\), antes de
clasificar trayectorias. La distancia puntual será
\begin{equation}
 d_S(x,y)=\lVert S^{-1}(x-y)\rVert_2.
 \label{camp:eq-scaled-distance}
\end{equation}
En presencia de un grupo finito de simetrías \(G\) que actúe por isometrías
de \(d_S\), la distancia entre órbitas de simetría será
\begin{equation}
 d_{S,G}(x,y)=\min_{g\in G}d_S(x,g y).
 \label{camp:eq-quotient-distance}
\end{equation}
La condición isométrica asegura que esta expresión define una métrica en el
cociente por \(G\). Si se cambian representantes a \(g_0x\) y \(h_0y\),
\[
 \min_{g\in G}d_S(g_0x,gh_0y)
 =\min_{g\in G}d_S(x,g_0^{-1}gh_0y)=d_{S,G}(x,y),
\]
pues \(g\mapsto g_0^{-1}gh_0\) permuta el grupo. La simetría se obtiene de
\(d_S(x,gy)=d_S(y,g^{-1}x)\). Si \(g,h\) realizan los mínimos para
\((x,y)\) y \((y,z)\), la desigualdad triangular es
\[
 d_{S,G}(x,z)\leq d_S(x,ghz)
 \leq d_S(x,gy)+d_S(gy,ghz)
 =d_{S,G}(x,y)+d_{S,G}(y,z).
\]
Finalmente, el mínimo es cero exactamente cuando \(x=gy\) para algún
\(g\), es decir, cuando ambos puntos representan la misma órbita de simetría.
La finitud de \(G\) garantiza que los mínimos se alcanzan.

Para una simetría afín \(g(x)=A_gx+b_g\), la hipótesis equivale a
\begin{equation}
 A_g^{\mathsf T}S^{-2}A_g=S^{-2},\qquad g\in G.
 \label{camp:eq-symmetry-isometry}
\end{equation}
Sin ella, \cref{camp:eq-quotient-distance} puede fallar incluso en la
simetría entre sus argumentos. Por ejemplo, si \(S=\operatorname{diag}(1,2)\)
y \(G\) intercambia las coordenadas, para \(x=(1,0)\), \(y=(0,2)\) se obtiene
\(d_{S,G}(x,y)=1\) y \(d_{S,G}(y,x)=1/2\). Una alternativa para un grupo
finito afín es promediar el tensor métrico:
\begin{equation}
 M_G=\frac1{|G|}\sum_{g\in G}A_g^{\mathsf T}S^{-2}A_g,
 \qquad d_{M_G}(x,y)=\sqrt{(x-y)^{\mathsf T}M_G(x-y)}.
 \label{camp:eq-symmetry-averaged-metric}
\end{equation}
El tensor es definido positivo y satisface
\(A_h^{\mathsf T}M_GA_h=M_G\) para todo \(h\in G\), por reindexación de la
suma. Con esta métrica, o con \cref{camp:eq-symmetry-isometry}, el mínimo
define una distancia entre clases y evita contar como candidatos independientes
dos semillas relacionadas por una simetría exacta. Las inversiones centrales y
los cambios de signo preservan cualquier \(S\) diagonal positivo.

En Caputo se distinguirán dos experimentos:

\begin{description}[leftmargin=3.7cm,style=nextline]
  \item[reinicio o corte reset]
  se prescribe \(x(t_0)=x_0\) y se inicia una memoria nueva en el mismo terminal
  inferior \(t_0\); dos semillas cercanas se comparan en esta sección puntual;

  \item[historia heredada]
  se conserva el perfil de memoria producido por la continuación o por la
  trayectoria anterior; el estado de la prueba es el perfil completo, no sólo
  su evaluación actual.
\end{description}

Las dos respuestas se registrarán por separado. Reiniciar desde el último
punto de una continuación no reproduce la continuación heredada. De igual
modo, una coincidencia en la proyección \(\mathbb R^n\) no implica proximidad
en el espacio funcional. Una afirmación sobre historias se formulará en un
espacio admisible \(X_{\mathrm{adm}}\subset\mathfrak C\), positivamente
invariante, con topología relativa y semiflujo bien definido y continuo, como
se precisa en \cref{sec:geometria-topologia-caputo}. La familia admisible y
su invariancia deben justificarse; no se presupone atracción ni ocultedad en
todo el espacio de términos libres \(\mathfrak C\). Cuando sea necesario
comparar perfiles se podrá usar una norma ponderada previamente fijada, por ejemplo
\begin{equation}
 d_{\eta,S}(m_1,m_2)=
 \sup_{\theta\geq0}e^{-\eta\theta}
 \lVert S^{-1}(m_1(\theta)-m_2(\theta))\rVert_2,
 \qquad \eta>0.
 \label{camp:eq-history-distance}
\end{equation}
Esta norma sólo es finita en la clase de perfiles con el crecimiento
ponderado requerido; no define automáticamente la topología compacto-abierta
usada para el semiflujo. Si se adopta como topología de trabajo, deben
verificarse en ella continuidad e invariancia de \(X_{\mathrm{adm}}\).
El valor de \(\eta\), la ventana en que se discretiza el perfil y el error de
truncamiento formarán parte del protocolo y se fijarán antes de evaluar el
resultado.

Para campos por partes se conservará, además, la convención sobre los planos
de conmutación, la dirección de cruce y el tratamiento de un posible contacto
tangencial. Una corrida que cambia de región sin localizar el evento no puede
servir para inferir la topología de una cuenca.

\subsection{Hipótesis falsables por familia de casos}
\label{camp:hipotesis-casos}

Las hipótesis siguientes no anticipan una respuesta afirmativa. Su función es
decidir qué observación refutaría la utilidad de cada localizador y qué método
complementario debe entrar después.

\begin{longtable}{@{}p{0.17\linewidth}p{0.35\linewidth}p{0.39\linewidth}@{}}
\caption{Hipótesis geométrico--topológicas que guían la campaña.}
\label{camp:tab-hypotheses}\\
\toprule
Caso & Hipótesis de localización & Resultado que la debilita o refuta\\
\midrule
\endfirsthead
\toprule
Caso & Hipótesis de localización & Resultado que la debilita o refuta\\
\midrule
\endhead
Chua PWL, \(q=1\)
& La continuación armónica y las intersecciones con \(x=\pm1\) aportan semillas
  a lados distintos de separatrices; la ausencia de PP regionales hace de este
  caso un control de complementariedad.
& Todas las semillas válidas llegan al mismo destino, escapan o la clasificación
  cambia al resolver exactamente los eventos de conmutación.\\

Chua PWL, Caputo
& Las mismas superficies de conmutación organizan la proyección, pero la
  separación de destinos depende de la historia. La continuación heredada y
  el reinicio desde su punto final pueden pertenecer a clases distintas.
& La diferencia desaparece bajo refinamiento de paso y memoria, o sólo aparece
  cuando se trunca la historia.\\

Chua arctangente, Wu
& Las semillas armónicas, el par FPP--A y la curva conectante son muestras
  geométricamente diferentes; su combinación puede revelar el mismo destino o
  dos dominios de atracción distintos.
& El supuesto destino recurrente no sobrevive a \(h/2\), al segundo integrador
  o a la ampliación del horizonte.\\

Chua arctangente, \texttt{c590}
& La gran norma del campo convierte el par FPP--A en un control de viabilidad
  dinámica, que se compara con las rutas armónica y de continuación.
& La acotación del FPP--A aparece sólo en horizontes cortos, o alguna ruta
  rebasa de manera reproducible la región acotada declarada.\\

Kalman--Fitts
& La inexistencia de PP por invertibilidad del Jacobiano no impide localizar
  fronteras mediante superficies críticas, retorno y grafo exterior. El caso
  mide cuánto se perdería al usar PP como localizador único.
& Las componentes recurrentes del grafo desaparecen al refinar cajas o son
  artefactos de una conmutación mal resuelta.\\

MAVPD, \(\xi=3.1\)
& Los PP conocidos funcionan como control de acceso a los ciclos interiores;
  las semillas de continuación y de borde son necesarias para estudiar el
  ciclo exterior y su separatriz.
& Los clasificadores no recuperan los destinos de referencia o el borde cambia
  de forma no convergente con tolerancia y horizonte.\\

MAVPD, \(\xi=3.5\)
& Al cambiar \(\xi\), la rama algebraica de PP persiste pero puede cambiar su
  destino. La campaña debe reproducir el papel de este parámetro como control
  de accesibilidad antes de discutir ocultedad.
& La asignación de cuenca contradice entre integradores o sólo se obtiene con
  una selección posterior de semillas.\\

MAVPD, \(\xi=2.85\)
& La continuación que llega al candidato caótico de referencia y el banco PP
  deben delimitar destinos distintos; el seguimiento de borde puede producir
  una órbita de borde susceptible de análisis de retorno.
& Los indicadores de caos, la recurrencia o la asignación de destino dejan de
  coincidir al refinar; entonces el objeto queda como transitorio pendiente.\\

PLL lead--lag
& La geometría correcta es cilíndrica. Los PP por vuelta, el retorno escalar y
  el borde entre régimen síncrono y corredor de fase se complementan.
& La separación sólo existe en la coordenada angular desenrollada o cambia al
  aplicar la métrica del cilindro.\\

Lorenz generalizado, \(q=0.995\)
& El par FPP--A y las semillas armónicas reducidas por la simetría
  \((x,y,z)\mapsto(-x,-y,z)\) pueden explorar cuencas distintas de las de los
  equilibrios.
& La nube de referencia resulta transitoria o cambia con la memoria o el
  horizonte; la cota de permanencia de \cref{prop:geofo-lorenz-boundedness}
  no decide ese destino.\\

Rabinovich--Fabrikant, \(q=0.998\)
& Los ocho FPP--A, reducidos a cuatro representantes por simetría, cubren
  regiones que una sola semilla armónica no visita; la prueba primaria es
  separar atracción, equilibrio y escape.
& Las semillas abandonan la región de prueba o las supuestas componentes se
  fusionan cuando aumenta la resolución.\\

Muñoz--Pacheco, \(q=0.97\) y \(q=0.996\)
& Las cuatro semillas FPP--A y la conmutación \(x=0\) permiten estudiar
  coexistencia y accesibilidad en ambos órdenes.
& El evento \(x=0\) no converge, la recurrencia desaparece o los dos órdenes
  sólo parecen distintos por la ventana temporal. En ausencia de equilibrios,
  el análisis se formula en términos de coexistencia, aislamiento y dependencia
  de la historia.\\
\bottomrule
\end{longtable}

\subsection{Banco unido de semillas y reducción por simetría}
\label{camp:banco-semillas}

La comparación asignará cuotas homogéneas a cada método y construirá la unión
etiquetada
\begin{equation}
 \mathcal S_0=
 \mathcal S_{\mathrm{DF}}
 \cup\mathcal S_{\mathrm{Mach}}
 \cup\mathcal S_{\mathrm{cont}}
 \cup\mathcal S_{\mathrm{PP/FPP}}
 \cup\mathcal S_{\mathrm{crit}}
 \cup\mathcal S_{\mathrm{conn}}
 \cup\mathcal S_{\mathrm{eig}}
 \cup\mathcal S_{\mathrm{edge}}.
 \label{camp:eq-seed-union}
\end{equation}
Cada semilla conservará su procedencia, parámetros, fase, amplitud, región del
campo, historia inicial y transformación de coordenadas. La deduplicación se
aplicará después de \cref{camp:eq-quotient-distance} y conservará la procedencia
de todos los métodos que produzcan representantes coincidentes.

\paragraph{Semillas armónicas.}
\(\mathcal S_{\mathrm{DF}}\) contendrá las soluciones de balance fundamental y,
cuando el residuo lo exija, las fases del balance multiharmónico. La ruta de
Machado \(\mathcal S_{\mathrm{Mach}}\) se usará cuando la aproximación racional
fraccionaria y su banda de frecuencia estén declaradas. Una raíz de la
ecuación de frecuencia que queda fuera de esa banda se conservará como control
rechazado y quedará excluida del banco activo. Las fases \(0,\pi/2,\pi,3\pi/2\) se
probarán cuando el balance no fija una fase única.

\paragraph{Continuación causal.}
\(\mathcal S_{\mathrm{cont}}\) incluirá estados finales de ramas de parámetro y,
para Caputo, el perfil heredado completo. Se guardarán dos hijos diferentes:
uno que conserva la historia y otro que reinicia en el mismo punto final. Esta
pareja cuantifica el efecto de memoria sin confundirlo con el efecto de la
semilla puntual.

\paragraph{Puntos perpetuos y FPP--A.}
\(\mathcal S_{\mathrm{PP/FPP}}\) contendrá todas las raíces reales no
estacionarias de \(J_f f=0\) admitidas por el dominio. En la ruta Caputo se
interpretarán únicamente como cancelación del segundo coeficiente de
Picard--Caputo, según \cref{sec:perpetual-points}. Se describirán como semillas
FPP--A, distintas de puntos de aceleración fraccionaria nula. Chua PWL y
Kalman--Fitts son controles sin PP; su banco contendrá sólo raíces que satisfagan
las ecuaciones.

\paragraph{Superficies críticas y curvas conectantes.}
Para cada componente se muestrearán
\begin{equation}
 \mathcal V_i=\{x:f_i(x)=0\},
 \qquad
 \mathcal A_i=\{x:(J_f(x)f(x))_i=0\},
 \label{camp:eq-critical-surfaces}
\end{equation}
y la variedad de alineación
\begin{equation}
 \mathcal C=
 \left\{x:f(x)\neq0,\
 \operatorname{rank}[f(x),J_f(x)f(x)]\leq1\right\}.
 \label{camp:eq-connecting-set}
\end{equation}
En tres dimensiones se verificará la ecuación equivalente
\(f\times(J_f f)=0\); en dimensión general se anularán todos los menores
\(2\times2\). Los PP son el subconjunto de \(\mathcal C\) con \(J_f f=0\), de
modo que las curvas conectantes amplían el banco sin invalidar el control
negativo de PP.

Para evitar aceptar una raíz espuria se registrarán dos residuos escalados:
\begin{align}
 r_{\mathrm{PP}}(x)
 &=\frac{\lVert J_f(x)f(x)\rVert_2}
        {1+\lVert J_f(x)\rVert_2\lVert f(x)\rVert_2},
 \label{camp:eq-pp-residual}\\
 r_{\mathrm{conn}}(x)
 &=\frac{\lVert \Pi_{f(x)}^{\perp}J_f(x)f(x)\rVert_2}
        {1+\lVert J_f(x)f(x)\rVert_2},
 \qquad
 \Pi_f^{\perp}=I-\frac{ff^{\mathsf T}}{f^{\mathsf T}f}.
 \label{camp:eq-connecting-residual}
\end{align}
En los campos por partes estas ecuaciones se resolverán región por región y se
comprobarán aparte los límites laterales de cada superficie de conmutación.
Para \(q<1\), \(\mathcal A_i\) y \(\mathcal C\) se interpretan como conjuntos
del desarrollo local de Picard--Caputo en la variable \(s=(t-t_0)^q\); no se
les atribuye el significado de aceleración física puntual.

\paragraph{Indicador KCC de priorización.}
Cuando se adopte una elevación de segundo orden explícita, se evaluará el
tensor de desviación \(P(x)\) de esa convención y el indicador
\begin{equation}
 \kappa_{\mathrm{KCC}}(x)=
 \max\{\Re\lambda:\lambda\in\operatorname{spec}P(x)\}.
 \label{camp:eq-kcc-score}
\end{equation}
Valores grandes y positivos priorizan regiones de separación transversal;
valores negativos se retienen como controles de regiones contractivas. Este
indicador sólo ordena semillas. Depende de la elevación, la métrica y las
coordenadas elegidas, y no demuestra caos, atracción ni ocultedad. Para evitar
sesgo, cada estrato conservará candidatos de ambos extremos y de la mediana de
\(\kappa_{\mathrm{KCC}}\).

\paragraph{Semillas espectrales y de borde.}
\(\mathcal S_{\mathrm{eig}}\) contendrá perturbaciones de cada equilibrio en
direcciones propias reales y en los planos reales de pares complejos. Si dos
semillas ya clasificadas tienen destinos diferentes, sus combinaciones
admisibles generan \(\mathcal S_{\mathrm{edge}}\). Estas últimas no se
construyen hasta que ambos extremos hayan pasado el control de convergencia.

\subsection{Las nueve etapas TG0--TG8}
\label{camp:etapas}

\subsubsection*{TG0: fijar el objeto y el protocolo}

Se registran \cref{camp:eq-ode-contract} o
\cref{camp:eq-caputo-contract}, terminal inferior, órdenes, parámetros,
escalas, simetrías, regiones no suaves, dominio de prueba y definición del
destino. También se fija si la afirmación buscada concierne al espacio de
estados, a la sección de reinicio o a un espacio funcional admisible
\(X_{\mathrm{adm}}\) con topología relativa. En el último caso se registran
también la familia de perfiles, su invariancia y la noción de atracción
empleada. Estos datos definen la entrada mínima de la campaña.

\subsubsection*{TG1: verificación algebraica antes de integrar}

Se calculan equilibrios, Jacobianos, espectros, criterio de estabilidad
correspondiente, simetrías, divergencia, planos de conmutación, funciones de
transferencia y residuos armónicos. Después se resuelven
\cref{camp:eq-critical-surfaces,camp:eq-connecting-set}. Las raíces se
verifican por sustitución a precisión mayor que la usada en la integración.
Esta etapa produce candidatos algebraicos; no produce atractores.

\subsubsection*{TG2: permanencia, escape y dominio computacional}

Antes de clasificar una nube como recurrente se busca una región absorbente o
una desigualdad de Lyapunov. Si no se dispone de ella, se declara una caja
computacional \(D\), un radio de escape \(R_{\mathrm{esc}}\), una cota sobre el
campo y una regla de ampliación. Una trayectoria se etiqueta
\emph{escape observado} sólo si cruza \(R_{\mathrm{esc}}\) con crecimiento
persistente y el evento se reproduce al refinar; de otro modo se etiqueta
\emph{indeterminada}. Para Lorenz generalizado, la cota de
\cref{prop:geofo-lorenz-boundedness} cierra existencia global y permanencia
para los parámetros declarados y los PVI de reinicio. La clasificación del
destino sigue pendiente. RF requiere todavía este criterio de permanencia
antes de TG5.

\subsubsection*{TG3: construir, estratificar y deduplicar semillas}

Se forma \(\mathcal S_0\) con cuotas por método, componente conexa, simetría y
escala. En el nivel inicial se tomarán todas las soluciones aisladas exactas,
al menos \(32\) puntos cuasiuniformes por componente accesible de superficie
crítica y \(64\) puntos por rama conectante. Los niveles siguientes elevan
esas cuotas a \(128/256\) y refinan sólo las celdas donde cambie el destino.
Ningún método recibe más presupuesto por haber generado primero una figura
atractiva.

\subsubsection*{TG4: transportar causalmente cada semilla}

En ODE se integra cada condición inicial y se conserva el tiempo de cada cruce
o retorno. En continuación paramétrica se registra el programa
\(\mu(t)\) y el tiempo de permanencia en cada escalón. En Caputo se realizan,
cuando tengan sentido, las dos ramas siguientes:
\begin{equation}
 \text{misma proyección final}\quad\Longrightarrow\quad
 \begin{cases}
   \text{continuación con historia heredada},\\
   \text{PVI nuevo con memoria reiniciada}.
 \end{cases}
 \label{camp:eq-history-pair}
\end{equation}
La memoria heredada se reconstruye desde el terminal inferior; un PVI Caputo
nuevo desde el último punto se registra como reinicio. En campos PWL cada lado del
evento se evalúa con su fórmula declarada y se rechaza una corrida si la
secuencia de regiones cambia sin converger con \(h\).

\subsubsection*{TG5: convergencia dinámica y clasificación de destinos}

Se aplican los presupuestos de \cref{camp:tab-budgets} y los observables de
\cref{camp:observables}. La concordancia punto a punto sólo se exige durante
una ventana corta, pues dos trayectorias caóticas pueden separarse aun cuando
representen el mismo invariante. A tiempos largos se comparan clasificador,
medidas empíricas, retornos y exponentes. Una etiqueta dinámica se acepta
cuando es estable frente al paso, el integrador y el horizonte declarados.

\ifHAFOBasinMaterial
\subsubsection*{TG6: geometría de cuenca y seguimiento de borde}

Se sondean vecindades de equilibrios, tubos alrededor de curvas conectantes y
celdas donde dos destinos vecinos difieren. Se refina la frontera por
bisección y continuación de borde, como se detalla en
\cref{camp:edge-tracking}. Esta etapa entrega una frontera numérica dependiente
de resolución; no entrega aún una variedad estable exacta.
\fi

\subsubsection*{TG7: grafo exterior, bloque aislante e índice}

Sólo las regiones persistentes pasan al método orientado a conjuntos. Se
construye una aplicación multivaluada exterior y se buscan componentes
fuertemente conexas, pares índice y bloques aislantes. Una integración de
centros de caja aporta una exploración preliminar; las imágenes se encierran con aritmética
intervalar o con cotas validadas. En Caputo la certificación se formula en un
espacio admisible de memoria o en una elevación finita identificada como
aproximación. Las cotas finitas de error de núcleo y solución no transportan
por sí solas el índice al sistema Caputo; se requiere el argumento de
continuación descrito en \cref{camp:caputo-lift}.

\subsubsection*{TG8: evaluar el alcance del resultado}

Cada caso termina en una de cuatro decisiones: \emph{concluir la rama},
\emph{reformular}, \emph{conservar como candidato} o \emph{ampliar el nivel de evidencia}. La decisión depende de los criterios de \cref{camp:evidence-levels} y se aplica
a cada rama del banco de semillas por separado.

\subsection{Presupuestos escalonados de integración, horizonte y memoria}
\label{camp:budgets}

Sea \(\tau_*>0\) una escala temporal fijada antes de integrar. Entre las
escalas disponibles pueden utilizarse el tiempo de cruce de la no linealidad,
el periodo asociado a una frecuencia de oscilación identificada y una escala
espectral. Para un modo \({}^{\mathrm C}D^q x=\lambda x\), el argumento de la
solución \(E_q(\lambda t^q)\) es adimensional; por tanto, \(\lambda\) tiene
unidades de tiempo a la potencia \(-q\) y, si \(\lambda\ne0\), su escala es
\(|\lambda|^{-1/q}\). La elección del menor tiempo espectral da
\[
 \tau_{\mathrm{esp}}=
 \left(\max_{\lambda\ne0}|\lambda|\right)^{-1/q}.
\]
Para \(q=1\) se recupera el inverso de la magnitud espectral. Un autovalor
complejo fraccionario no implica por sí solo una respuesta periódica. Se
adoptará la menor de las escalas justificadas para el caso; si ya existe una
escala validada más estricta, prevalecerá esta última.

La transformación del tiempo también modifica el campo. Para mostrar el
factor preciso, sean \(t=t_0+\tau_*\theta\) y
\(Y(\theta)=S^{-1}X(t_0+\tau_*\theta)\), con la matriz de escalas de
\cref{camp:eq-scaled-distance}. En la ecuación de Volterra de
\({}^{\mathrm C}D_{t_0}^qX=F(t,X)\), el cambio
\(s=t_0+\tau_*u\) produce
\[
 (t-s)^{q-1}\,\dd s
 =\tau_*^{q-1}(\theta-u)^{q-1}\tau_*\,\dd u
 =\tau_*^q(\theta-u)^{q-1}\,\dd u.
\]
Multiplicar la ecuación por \(S^{-1}\) da entonces
\begin{equation}
 \begin{aligned}
 Y(\theta)={}&S^{-1}X_0\\
 &+\frac{\tau_*^q}{\Gamma(q)}\int_0^\theta(\theta-u)^{q-1}
       S^{-1}F(t_0+\tau_*u,SY(u))\,\dd u.
 \end{aligned}
 \label{camp:eq-time-scaled-volterra}
\end{equation}
Por la equivalencia integral, el problema adimensional es
\begin{equation}
 {}^{\mathrm C}D_0^qY(\theta)
 =\tau_*^q S^{-1}F(t_0+\tau_*\theta,SY(\theta)),
 \qquad Y(0)=S^{-1}X_0.
 \label{camp:eq-time-scaled-caputo}
\end{equation}
El factor es \(\tau_*^q\), y sólo se reduce a \(\tau_*\) en el caso entero.
Un paso físico \(h\) y una duración física \(T\), medida desde \(t_0\),
corresponden a \(\widehat h=h/\tau_*\) y
\(\widehat T=T/\tau_*\). Así, los presupuestos siguientes comparan resolución
temporal relativa a cada sistema, conservando \(N=T/h=\widehat T/\widehat h\).

\begin{table}[H]
\centering
\caption{Presupuesto creciente. \(N=\widehat T/\widehat h\) es el número
nominal de pasos; B3 sólo se aplica a candidatos que superaron B2.}
\label{camp:tab-budgets}
\begin{tabularx}{\linewidth}{@{}lccccY@{}}
\toprule
Nivel & \(\widehat h\) & \(\widehat T\) & Fracción transitoria & Semillas & Función\\
\midrule
B0 & \(2\times10^{-2}\) & \(50\)  & \(1/2\) & banco mínimo &
criba de errores, eventos y escapes evidentes\\
B1 & \(10^{-2}\) & \(100\) & \(1/2\) & banco completo inicial &
primera asignación de destino\\
B2 & \(5\times10^{-3}\) & \(200\) & \(3/5\) & semillas ambiguas y frontera &
convergencia de clasificación y observables\\
B3 & \(2.5\times10^{-3}\) & \(400\) & \(2/3\) & candidatos con evidencia B2 &
retornos largos, borde y preparación topológica\\
\bottomrule
\end{tabularx}
\end{table}

En ODE el integrador principal será adaptativo de orden alto, con tolerancias
\((\mathrm{rtol},\mathrm{atol})=(10^{-10},10^{-12})\), paso máximo adimensional
no mayor que el \(\widehat h\) del nivel y localización explícita de eventos. El control
independiente será un método implícito apropiado para rigidez o un RK de paso
fijo con \(\widehat h/2\); no compartirá interpolador ni detector de eventos.

En Caputo el integrador de referencia de memoria completa será ABM--PECE con suma desde \(t_0\).
Se compararán \(h,h/2,h/4\) en ventanas comunes. El segundo solucionador será
una cuadratura de convolución BDF2 o EFORK--3 sólo después de comprobar que
implementa el mismo operador, terminal inferior, orden y condición inicial.
Una coincidencia con un método que resuelve otra derivada fraccionaria no es
validación cruzada.
El refinamiento conservará el tratamiento del arranque y de las interfaces
temporales. Las condiciones que permiten estimar los defectos de cuadratura,
y su distinción respecto de un error global, se desarrollan en
\cref{subsec:numerical-startup-consistency}. Si se usa continuación por
saltos, se aplica \cref{subsec:numerical-continuation-interfaces}; si se
compara con un actualizador local reiniciado, se conserva la distinción
analítica de \cref{subsec:numerical-adm-history}.

El costo cuadrático del integrador de referencia limita su uso en B3.
Ese nivel empleará un algoritmo acelerado sólo después de validarlo en B0--B2,
y conservará ventanas y semillas testigo calculadas con la suma histórica
directa. Cualquier reducción de la longitud de memoria se declara y valida por
separado.

La comparación larga de dos trayectorias caóticas usará ventanas estadísticas.
En la ventana corta \(0\leq t-t_0\leq10\tau_*\) se exigirá
\begin{equation}
 \max_t d_S(x_h(t),x_{h/2}(t))\leq10^{-3}.
 \label{camp:eq-short-error}
\end{equation}
Después se exigirá igualdad del destino y diferencias relativas no mayores de
\(5\%\) en observables robustos. Si un observable cruza cero, se usará una
tolerancia absoluta predeclarada en vez de un cociente inestable.

\subsection{Elevación de memoria para Caputo y su control de error}
\label{camp:caputo-lift}

Para hacer viable el grafo exterior se podrá aproximar el kernel por una suma
de exponenciales,
\begin{equation}
 K_q(t)=\frac{t^{q-1}}{\Gamma(q)}
 \approx K_q^{(m)}(t)=\sum_{\ell=1}^{m}w_\ell e^{-\eta_\ell t},
 \qquad w_\ell,\eta_\ell>0,
 \label{camp:eq-soe-kernel}
\end{equation}
La positividad admite una representación integral exacta. En efecto,
\[
 \int_0^\infty\eta^{-q}\e^{-\eta t}\,\dd\eta
 =t^{q-1}\int_0^\infty u^{-q}\e^{-u}\,\dd u
 =t^{q-1}\Gamma(1-q).
\]
La identidad \(\Gamma(q)\Gamma(1-q)=\pi/\sin(\pi q)\) da
\[
 K_q(t)=\frac{\sin(\pi q)}{\pi}
        \int_0^\infty\eta^{-q}\e^{-\eta t}\,\dd\eta,
 \qquad t>0.
\]
Una cuadratura positiva en la variable \(\eta\) motiva las tasas y pesos
positivos de \cref{camp:eq-soe-kernel}. Para la suma finita se definen
\[
 z_\ell(t)=\int_0^t\e^{-\eta_\ell(t-s)}f(x(s))\,\dd s.
\]
La regla de Leibniz produce
\(\dot z_\ell=f(x(t))-\eta_\ell z_\ell\), con \(z_\ell(0)=0\). Sustituir
la suma en la ecuación de Volterra permite elevar el problema de núcleo
aproximado mediante variables de memoria,
escribiendo \(t_0=0\) sin pérdida de generalidad:
\begin{equation}
 \dot z_\ell=-\eta_\ell z_\ell+f(x),
 \qquad
 x=x_0+\sum_{\ell=1}^{m}w_\ell z_\ell,
 \qquad z_\ell(0)=0.
 \label{camp:eq-soe-lift}
\end{equation}
La elevación es exacta para el núcleo aproximado, no para
\cref{camp:eq-caputo-contract}; \(x_0\) sigue siendo un parámetro de esta
familia de sistemas. Para \(0<q<1\), el núcleo exacto es singular en cero y
una suma finita de exponenciales es acotada allí. Por ello no existe una cota
uniforme finita del error en \((0,T]\). Si se aproxima uniformemente sólo en
\([\delta,T]\), con \(\delta=h_{\min}>0\), el tramo local debe tratarse aparte.
Con corrección local exacta, la ecuación de reconstrucción pasa a ser
\begin{align}
 x_m(t)&=x_0+\sum_{\ell=1}^{m}w_\ell z_\ell(t)+R_{\mathrm{loc}}(t),\\
 R_{\mathrm{loc}}(t)
 &=\int_{\max(0,t-\delta)}^t
   [K_q(t-s)-K_q^{(m)}(t-s)]f(x_m(s))\,\mathrm ds.
 \label{camp:eq-soe-local-correction}
\end{align}
Se conserva \(\dot z_\ell=-\eta_\ell z_\ell+f(x_m)\). Esta corrección
requiere la historia local; el algoritmo resultante ya no es solamente la ODE
finita de \cref{camp:eq-soe-lift}. Equivalentemente, su núcleo
\(\widetilde K_m\) coincide con \(K_q\) en \((0,\delta)\) y con
\(K_q^{(m)}\) en \([\delta,T]\). Si se usa otra corrección, se declara el
núcleo efectivo o el defecto integral correspondiente.

El control relevante para la ecuación integral es
\begin{equation}
 \varepsilon_K^{(m)}(T)=
 \int_0^{T}\left|K_q(t)-\widetilde K_m(t)\right|\,\mathrm dt
 \leq\varepsilon_{\mathrm{loc}}+(T-\delta)\varepsilon_\infty,
 \label{camp:eq-kernel-error}
\end{equation}
para \(T\geq\delta\), donde \(\varepsilon_{\mathrm{loc}}\) es el error
\(L^1(0,\delta)\) del núcleo efectivo y \(\varepsilon_\infty\) acota el error
en \([\delta,T]\). La corrección local exacta da
\(\varepsilon_{\mathrm{loc}}=0\). Sin corrección local, una cota explícita del tramo singular se obtiene por
la desigualdad triangular:
\[
 \int_0^\delta|K_q(t)-K_q^{(m)}(t)|\,\dd t
 \leq\frac{\delta^q}{\Gamma(q+1)}
       +\sum_{\ell=1}^m\frac{w_\ell}{\eta_\ell}
                         (1-\e^{-\eta_\ell\delta}).
\]
El primer término integra exactamente la singularidad y cada sumando integra
una exponencial. La cota \cref{camp:eq-kernel-error} se obtiene al sumar este
error local, o el de la corrección elegida, y
\(\int_\delta^T|K_q-\widetilde K_m|\leq(T-\delta)\varepsilon_\infty\).
El error de una cuadratura local o histórica
se registra, con unidades de estado, como defecto integral adicional
\(\varepsilon_{\mathrm{quad}}\), y no se confunde con el error del núcleo.

\begin{proposition}[Control del error en horizonte finito]
\label{camp:prop-soe-error-finite}
Sean \(x\) la solución Caputo y \(x_m\) una aproximación continua que comparte
\(x_0\) y satisface
\[
 x_m(t)=x_0+\int_0^t\widetilde K_m(t-s)f(x_m(s))\,\mathrm ds+r_m(t),
 \qquad \|r_m\|_{\infty,[0,T]}\leq\varepsilon_{\mathrm{quad}}.
\]
Supóngase que ambas curvas permanecen en una región donde \(f\) es
\(L\)-Lipschitz y \(\|f\|\leq M\). Entonces
\begin{equation}
 \sup_{0\leq t\leq T}\|x(t)-x_m(t)\|
 \leq\bigl[M\varepsilon_K^{(m)}(T)+\varepsilon_{\mathrm{quad}}\bigr]
       E_q(LT^q).
 \label{camp:eq-soe-solution-error}
\end{equation}
\end{proposition}

\begin{proof}
Al restar las ecuaciones y sumar y restar la convolución
\(K_q*f(x_m)\), se obtiene la identidad
\begin{align*}
 x(t)-x_m(t)
 &=\int_0^tK_q(t-s)[f(x(s))-f(x_m(s))]\,\dd s\\
 &\quad+\int_0^t[K_q(t-s)-\widetilde K_m(t-s)]f(x_m(s))\,\dd s-r_m(t).
\end{align*}
La primera integral se acota mediante la constante de Lipschitz, la segunda
mediante \(M\varepsilon_K^{(m)}(T)\), y el último término mediante
\(\varepsilon_{\mathrm{quad}}\). Con
\(u(t)=\|x(t)-x_m(t)\|\) y
\(A=M\varepsilon_K^{(m)}(T)+\varepsilon_{\mathrm{quad}}\),
\[
 u(t)\leq A+L(I^qu)(t),\qquad
 (I^qv)(t)=\frac1{\Gamma(q)}\int_0^t(t-s)^{q-1}v(s)\,\dd s.
\]
Para iterar esta desigualdad se usa la identidad de convolución
\begin{equation}
 \frac1{\Gamma(a)\Gamma(b)}
 \int_0^t(t-s)^{a-1}s^{b-1}\,\dd s
 =\frac{t^{a+b-1}}{\Gamma(a+b)},\qquad a,b>0.
 \label{camp:eq-fractional-convolution}
\end{equation}
Se demuestra con \(s=tv\): la integral restante es
\(B(a,b)=\Gamma(a)\Gamma(b)/\Gamma(a+b)\). Aplicada a integrales iteradas,
con Fubini, implica \(I^aI^b=I^{a+b}\) e
\(I^{jq}1=t^{jq}/\Gamma(jq+1)\). Tras \(N\) sustituciones,
\[
 u(t)\leq A\sum_{j=0}^{N-1}\frac{L^jt^{jq}}{\Gamma(jq+1)}
             +L^N(I^{Nq}u)(t).
\]
La continuidad de ambas curvas en \([0,T]\) da \(\|u\|_\infty<\infty\) y
\[
 0\leq L^N(I^{Nq}u)(t)
 \leq\|u\|_\infty\frac{(LT^q)^N}{\Gamma(Nq+1)}\longrightarrow0.
\]
La última convergencia sigue del crecimiento de la función gamma. Al pasar
al límite, la serie es \(E_q(Lt^q)\); como sus coeficientes son no negativos,
\(E_q(Lt^q)\leq E_q(LT^q)\). Esto prueba
\cref{camp:eq-soe-solution-error}, en la forma de Gronwall fraccionaria de
\cite{Gomoyunov2018Convex}. La permanencia en la región debe justificarse
mediante una envolvente validada o un argumento de continuación para usar la
cota como certificado.
\end{proof}

También se medirá el error frente al integrador de referencia de memoria
completa en B0--B2; esa comparación numérica no sustituye las cotas del lema. Los
presupuestos serán \(m=8,16,32,64\). Se detendrá el refinamiento sólo cuando
dos valores consecutivos conserven destinos y métricas dentro de las
tolerancias de TG5, como criterio de estabilidad numérica en la ventana.

El límite asintótico requiere un análisis separado. Para \(f(x)=-ax\),
\(a>0\), la elevación sin corrección local de \cref{camp:eq-soe-lift} tiene
equilibrio proyectado
\begin{equation}
 x_m^*=\frac{x_0}{1+a\sum_{\ell=1}^{m}w_\ell/\eta_\ell},
 \label{camp:eq-soe-biased-equilibrium}
\end{equation}
En efecto, \(\dot z_\ell=0\) impone
\(z_\ell^*=-ax_m^*/\eta_\ell\); la reconstrucción da
\(x_m^*=x_0-ax_m^*\sum_\ell w_\ell/\eta_\ell\). Reunir los términos que
contienen \(x_m^*\) produce \cref{camp:eq-soe-biased-equilibrium}.
Para \(a=w_1=\eta_1=x_0=1\), se tiene \(x_m=1+z\) y
\(\dot z=-z-(1+z)=-1-2z\), con \(z(0)=0\). La solución es
\(z(t)=(\e^{-2t}-1)/2\) y, por tanto,
\(x_m(t)=(1+\e^{-2t})/2\), mientras la solución
Caputo \(x(t)=E_q(-t^q)\) tiende a cero. Una SOE finita con tasas positivas
puede alterar el equilibrio límite. Deben controlarse las bajas frecuencias,
la cola de memoria y el papel de \(x_0\) antes de comparar conjuntos
asintóticos o cuencas de diferentes reinicios; en particular, el orden de los
límites \(m\to\infty\) y \(T\to\infty\) necesita justificación.

Transportar un índice de Conley de la elevación al problema Caputo exige un
teorema de continuación aplicable. Su demostración deberá incluir:

\begin{enumerate}
  \item un espacio de fases común admisible, o inmersiones e identificaciones
        justificadas de los estados de memoria, en el que estén definidos el
        problema exacto, las aproximaciones y una homotopía continua;
  \item las hipótesis de existencia, continuidad y admisibilidad del
        semiflujo y de la homotopía que requiere el índice empleado;
  \item una vecindad aislante común durante toda la homotopía, con aislamiento
        uniforme y control de las contribuciones de memoria y de los errores
        intervalares; las cotas finitas de
        \cref{camp:eq-soe-solution-error} son una herramienta, no esta prueba;
  \item la correspondencia entre los conjuntos invariantes de los extremos y
        el teorema que conserva su índice. La persistencia observada al
        aumentar \(m\) y \(T\) será un control adicional de la construcción.
\end{enumerate}

Coincidencias en un número finito de valores de \(m\) y \(T\), o una cota
de error en una ventana, no prueban por sí solas esas hipótesis. Hasta cerrar
este argumento, el índice calculado pertenece al modelo elevado y el lift
sirve como acelerador e instrumento de exploración del problema Caputo.

\subsection{Observables y clasificador de destino}
\label{camp:observables}

Cada trayectoria producirá observables con unidades y ventanas explícitas:

\begin{enumerate}
  \item radio máximo, velocidad de crecimiento y distancia a la frontera de la
        región de prueba;
  \item distancia a cada equilibrio y norma de \(f(x)\);
  \item medias, covarianzas y cuantiles de cada coordenada después del
        transitorio;
  \item histograma o medida empírica sobre una sección transversal;
  \item tiempos de retorno, dispersión del retorno y número de cruces válidos;
  \item espectro de Lyapunov cuando la formulación variacional sea válida;
  \item prueba \(0--1\), entropía o complejidad como diagnósticos
        complementarios a las pruebas de convergencia;
  \item para Caputo, sensibilidad al reinicio, a la historia heredada y al
        número de modos de memoria.
\end{enumerate}

La clasificación no forzará a toda trayectoria a una categoría conocida. Se
usará el conjunto
\begin{equation}
 \mathfrak D=\{\text{equilibrio},\text{destino periódico ODE},
 \text{recurrente proyectado},\text{escape},\text{transitorio},
 \text{ambiguo}\}.
 \label{camp:eq-destination-set}
\end{equation}
En Caputo se evita llamar \emph{ciclo límite} a una nube casi periódica; se
informa como recurrencia proyectada salvo que exista un resultado teórico
compatible con la memoria. Un exponente máximo se clasificará como positivo
cuando converja con el paso, la ventana, el algoritmo QR y el modelo de memoria.

La recurrencia se cuantifica mediante dos conteos. Sea \(N_{\rm pares}\) el
número de pares temporales admisibles que satisfacen el criterio de retorno,
y \(N_{\rm origen}\) el número de muestras distintas desde las cuales se
observa al menos un retorno. Si \(N\) es el número de muestras evaluadas,
\[
 \nu_{\rm retorno}=\frac{N_{\rm origen}}{N}\in[0,1],
 \qquad
 m_{\rm retorno}=\frac{N_{\rm pares}}{N}.
\]
La primera magnitud mide la fracción de muestras con retorno; la segunda mide
su multiplicidad y puede superar uno. Ambas describen recurrencia dentro de
la ventana observada.

Para comparar dos nubes \(A\) y \(B\) se emplearán simultáneamente distancia
de centroides, diferencia de covarianzas y una distancia entre medidas
empíricas. Dos nubes se considerarán concordantes únicamente cuando los tres
indicadores cumplan las tolerancias declaradas; en otro caso el destino se
clasificará como \emph{ambiguo}.

\ifHAFOBasinMaterial
\subsection{Sondeo de cuencas y seguimiento de borde}
\label{camp:edge-tracking}

Alrededor de cada equilibrio \(e\) se fijará la escala
\(s_e=\max(1,\lVert S^{-1}e\rVert_2)\) y los radios
\begin{equation}
 r_j=s_e10^{-j},\qquad j=1,2,\ldots,6.
 \label{camp:eq-equilibrium-radii}
\end{equation}
En cada radio se usarán las direcciones propias admisibles y, como mínimo,
\(32\) direcciones cuasiuniformes adicionales. El resultado siempre se
reportará con su radio: cero contactos en una muestra finita no demuestra que
toda una vecindad esté fuera de la cuenca.

Dados dos datos \(a_0,b_0\) con destinos distintos, el borde se aproxima así:

\begin{enumerate}
  \item confirmar ambos destinos en B1 y B2;
  \item formar el punto medio en coordenadas escaladas y resolver desde el
        terminal inferior, sin empalmar trayectorias;
  \item sustituir el extremo que comparte el destino del punto medio;
  \item repetir hasta que \(d_S(a_k,b_k)<10^{-8}\) o hasta que el clasificador
        sea ambiguo en tres refinamientos consecutivos;
  \item integrar el par estrecho, detectar el intervalo en que permanece cerca
        del borde y volver a encuadrar antes de que caiga a un destino;
  \item registrar retornos de la trayectoria de borde y repetir con al menos
        tres pares iniciales no colineales.
\end{enumerate}

En el PLL la interpolación se hará en la carta cilíndrica que minimiza la
distancia angular, no a través del corte \(0=2\pi\). En Caputo reset se
interpolan datos iniciales en \(t_0\). Para historia heredada no se interpolan
arbitrariamente dos perfiles: se perturba una historia admisible mediante una
familia de inicialización declarada, o se reconstruyen ambos extremos desde
el mismo pasado. Cada punto medio se integra de nuevo con su memoria completa.

La geometría del borde se cuantificará mediante el exponente de incertidumbre
y la entropía de cuenca en resoluciones sucesivas. Una afirmación Wada sólo se
considerará si existen al menos tres destinos y un test de frontera común
estable al refinar; una imagen coloreada no es evidencia suficiente.
\fi

\subsection{Grafo exterior, aislamiento e índice de Conley}
\label{camp:set-oriented}

Sea \(D=\bigcup_i B_i\) una unión de cajas y \(\Phi_\tau\) la aplicación de
tiempo o el retorno a una sección. La campaña construirá una aplicación
multivaluada exterior
\begin{equation}
 \mathcal F(B_i)=
 \{B_j:B_j\cap\widehat\Phi_\tau(B_i)\neq\varnothing\},
 \qquad
 \Phi_\tau(B_i)\subseteq\widehat\Phi_\tau(B_i),
 \label{camp:eq-outer-map}
\end{equation}
donde la segunda inclusión debe estar validada. Las componentes fuertemente
conexas del grafo son candidatas recurrentes. Si \(N\) es una unión de cajas,
se exigirá
\begin{equation}
 \operatorname{Inv}(N)\subset\operatorname{int}N
 \label{camp:eq-isolating-block}
\end{equation}
antes de llamarlo bloque aislante. El conjunto de salida y el par índice se
construirán con la misma aplicación exterior; su homología dará el índice de
Conley computado.

El mallado inicial será \(16^3\) para los casos tridimensionales, pero sólo se
refinarán cajas visitadas, fronteras de cuenca y componentes recurrentes. En
Kalman--Fitts se preferirá una sección tridimensional a una malla uniforme en
cuatro dimensiones. En el PLL se mallará una carta del cilindro con
identificación periódica. Para la elevación Caputo se usarán representaciones
dispersas y sólo después del control de núcleo y solución de
\cref{camp:eq-kernel-error,camp:eq-soe-solution-error}. El transporte del índice
exige además la homotopía y el aislamiento común de \cref{camp:caputo-lift};
una malla de centros en dimensión \(n(m+1)\) no es una certificación.

La ausencia de camino en el grafo desde cajas que cubren vecindades completas
de los equilibrios hasta el bloque candidato será una afirmación relativa a
\(D\), \(\tau\) y la resolución. Para convertirla en exclusión rigurosa se
necesitan imágenes exteriores intervalares de todas esas cajas y una cobertura
de la vecindad, no sólo sondas puntuales. Aun entonces, el índice prueba
invariancia aislada, no por sí solo atracción, caos u ocultedad; estas
propiedades requieren evidencia adicional.

\subsection{Controles positivos, negativos y de consistencia}
\label{camp:controls}

Los controles se ejecutan con la misma infraestructura que los candidatos:

\begin{enumerate}
  \item una ODE lineal estable con solución exacta verifica integración,
        clasificación de equilibrio y grafo sin recurrencia espuria;
  \item la ecuación escalar de Caputo con solución de Mittag--Leffler verifica
        terminal inferior, memoria completa y convergencia \(h,h/2,h/4\);
  \item el límite \(q\uparrow1\) se compara con la ODE sin exigir igualdad a
        un \(q<1\) fijo;
  \item Chua PWL y Kalman--Fitts verifican que el módulo PP devuelva un conjunto
        vacío, no minimizadores falsos;
  \item MAVPD \(\xi=3.1\) usa sus destinos reproducidos como control del
        clasificador PP y MAVPD \(\xi=3.5\) como control de sensibilidad de
        accesibilidad;
  \item las semillas de referencia de Chua, PLL, Lorenz, RF y Muñoz calibran
        la asignación de destino, sin convertir la reproducción anterior en
        resultado de la nueva campaña;
  \item cada simetría se prueba con pares \(x_0,gx_0\): las trayectorias deben
        conservar la relación prevista dentro del error;
  \item en Caputo se compara deliberadamente una historia heredada con su
        reinicio puntual. Que difieran es una respuesta física posible; que el
        programa las trate como idénticas es un fallo;
  \item se ejecutan semillas alejadas de toda raíz armónica o conectante para
        estimar la tasa de falsos hallazgos debidos a una clasificación laxa.
\end{enumerate}

\subsection{Criterios de cierre, reformulación y conservación}
\label{camp:decision-rules}

\paragraph{Concluir una rama.}
La rama concluye si el sistema no satisface las hipótesis del protocolo, si el residuo algebraico no baja
con precisión, si el escape se reproduce en B1--B2, si un evento PWL no
converge o si el segundo integrador resuelve un operador fraccionario distinto.
Las demás ramas de \cref{camp:eq-seed-union} se evalúan de manera independiente.

\paragraph{Reformular.}
Se reformula si el destino cambia con \(h\), si el horizonte sólo muestra un
transitorio, si el clasificador depende de una escala mal condicionada, si el
lift cambia de topología al aumentar \(m\), o si el borde desaparece al usar
la métrica correcta. La reformulación puede cambiar escala, sección o dominio; su evaluación debe
distinguir el efecto de cada cambio.

\paragraph{Conservar como candidato.}
Se conserva cuando existe acotación finita, recurrencia y concordancia de
clasificador entre B1--B2 y dos solucionadores, pero todavía no hay bloque
aislante ni exclusión rigurosa de cuencas de equilibrio. En este nivel la
clasificación es \emph{candidato}; las categorías \emph{atractor oculto
demostrado} y \emph{caos certificado} requieren la evidencia adicional de los
niveles posteriores.

\paragraph{Ampliar el nivel de evidencia.}
El nivel dinámico siguiente requiere B3, observables convergentes,
perturbaciones que regresan y una frontera reproducible; el nivel topológico
requiere además TG7. En Caputo el resultado especificará si concierne a la
sección de reinicio, al espacio admisible \(X_{\mathrm{adm}}\) con su topología
relativa o a la elevación aproximada. La transferencia de una propiedad
asintótica de la elevación al problema Caputo requiere el argumento específico
de \cref{camp:caputo-lift}; una cota de error finito no basta.

\subsection{Escala de evidencia EV--TG0--EV--TG7}
\label{camp:evidence-levels}

La notación \(\mathrm{EV}\text{-}\mathrm{TG}k\) indica el alcance de la
evidencia. Las etapas TG0--TG8 designan las operaciones del protocolo.

\begin{table}[H]
\centering
\caption{Niveles de evidencia exigibles a cada afirmación.}
\label{camp:tab-evidence-levels}
\begin{tabularx}{\linewidth}{@{}lY@{}}
\toprule
Nivel & Contenido mínimo\\
\midrule
EV--TG0 & Modelo, espacio, derivada, parámetros, métrica y afirmación objetivo
fijados; no hay resultado dinámico.\\
EV--TG1 & Equilibrios, estabilidad, simetría, residuos armónicos, superficies,
curvas y semillas verificados algebraicamente.\\
EV--TG2 & Trayectoria causal finita acotada en una ventana y dominio declarados;
no prueba atracción ni recurrencia asintótica.\\
EV--TG3 & Destino e indicadores robustos a paso, horizonte, tolerancia y segundo
solucionador; en Caputo, también a memoria.\\
EV--TG4 & Familia de semillas, sondeo de vecindades y borde adaptativo
reproducibles; la cuenca sigue siendo una aproximación finita.\\
EV--TG5 & Retorno o aplicación exterior validados, bloque aislante e índice
persistentes bajo refinamiento.\\
EV--TG6 & Atracción y exclusión certificada de vecindades completas de todos los
equilibrios dentro del dominio declarado: ocultedad en ODE o en la sección
reset especificada.\\
EV--TG7 & Formulación y exclusión de vecindades abiertas en la topología relativa
de un espacio admisible invariante \(X_{\mathrm{adm}}\), más certificación
independiente del carácter caótico si se afirma; nivel funcional condicionado
a la elección y validación de ese espacio.\\
\bottomrule
\end{tabularx}
\end{table}

Una muestra finita alrededor de equilibrios puede alcanzar como máximo
EV--TG4. Para EV--TG6 hace falta cubrir una vecindad completa mediante
enclosures validados o un argumento analítico. En Muñoz--Pacheco, donde no hay
equilibrios, EV--TG6 se formula mediante coexistencia, accesibilidad y frontera;
la definición basada en cuencas de equilibrio no discrimina ese caso.

\subsection{Evaluación por familia de sistemas}
\label{camp:case-execution}

\begin{longtable}{@{}p{0.16\linewidth}p{0.22\linewidth}p{0.25\linewidth}p{0.28\linewidth}@{}}
\caption{Orden de ejecución y criterio mínimo. La última columna distingue la
ejecución completada de los pasos que siguen pendientes.}
\label{camp:tab-case-execution}\\
\toprule
Caso & Semillas prioritarias & Geometría y controles & Salida ejecutada o pendiente\\
\midrule
\endfirsthead
\toprule
Caso & Semillas prioritarias & Geometría y controles & Salida ejecutada o pendiente\\
\midrule
\endhead
Chua PWL, \(q=1\)
& DF sesgada, continuación, \(x=\pm1\), conectantes; PP vacío.
& Simetría central y eventos laterales exactos.
& B0--B2: mapa de destinos y borde; B3: grafo exterior por regiones.\\

Chua PWL, fraccionario
& Machado/DF, continuación heredada y reset, superficies críticas.
& Misma conmutación; memoria completa y \(m=8,16,32,64\).
& B0--B2: diferencia heredada--reset; B3: persistencia del borde.\\

Wu, \(q=0.99\)
& Semillas publicadas, DF/Machado, par FPP--A y conectante.
& Simetría central; la periodicidad exacta no se presupone.
& Ejecutado B0--B2: FPP--A acotados sin destino común, EV--TG2; convergencia y
borde siguen pendientes.\\

\texttt{c590}, \(q=0.9999\)
& Continuación y armónicas primero; FPP--A como control de alta velocidad.
& Radio de escape y norma del campo predeclarados.
& B0--B2: destino o escape reproducible sin selección posterior.\\

Kalman--Fitts
& Retorno, superficies críticas, conectantes y semillas espectrales; PP vacío.
& Sección tridimensional del sistema de cuarto orden.
& B0--B2: destinos; B3: componentes recurrentes persistentes.\\

MAVPD, \(\xi=3.1\)
& PP para ciclos interiores; DF y continuación para rama exterior; edge seeds.
& Simetría central y destinos de referencia como controles.
& \ifHAFOBasinMaterial
Ejecutado B0--B2: destinos interiores/exterior y tercer destino en bracket,
EV--TG2; separatriz reproducible pendiente.
\else
Nivel máximo conservador EV--TG2.
\fi\\

MAVPD, \(\xi=3.5\)
& Mismo banco transportado con las cuotas predeclaradas.
& Control de accesibilidad y comparación con \(\xi=3.1\).
& B1--B2: diagrama que confirme o refute la separación.\\

MAVPD, \(\xi=2.85\)
& Continuación desde la rama, PP, sección y puntos de borde.
& Control caótico de referencia; dos cálculos variacionales.
& B1--B3: retornos, exponentes y bloque candidato convergentes.\\

PLL lead--lag
& PP por vuelta, semilla T1+A1, retorno escalar y borde.
& Cilindro, ángulo módulo \(2\pi\) y carta sin corte artificial.
& \ifHAFOBasinMaterial
Ejecutado B0--B2: $p_+\to E_{\rm focus}$, EV--TG3; $p_-$ y bracket sin
resolver, EV--TG4 pendiente.
\else
Nivel máximo conservador EV--TG3.
\fi\\

Lorenz generalizado, \(q=0.995\)
& Par FPP--A, armónicas multiharmónicas y semillas de equilibrio.
& Simetría \((-x,-y,z)\); cota analítica de TG2 para reinicios.
& Permanencia demostrada; B0--B2: destino; después, borde en el espacio
admisible declarado.\\

RF, \(q=0.998\)
& Cuatro representantes FPP--A, armónicas y semillas espectrales.
& Simetría \((-x,-y,z)\), cinco equilibrios y control de escape.
& TG2 y B0--B3: separación de destinos y bloque persistente.\\

Muñoz, \(q=0.97\)
& Cuatro FPP--A, \(x=0\), conectantes y banco armónico.
& Simetría \((-x,y,-z)\); no hay equilibrios.
& B0--B2: convergencia y mapa de coexistencia, sin rótulo oculto.\\

Muñoz, \(q=0.996\)
& Mismo banco y semillas de referencia, sin cambiar clasificador.
& Comparación de memoria y horizonte con \(q=0.97\).
& B1--B3: persistencia o desaparición de la coexistencia proyectada.\\
\bottomrule
\end{longtable}

\subsection{Orden práctico de la campaña}
\label{camp:execution-order}

B0--B2 cubre MAVPD $\xi=3.1$, PLL y Wu. La ampliación comienza con
controles analíticos y enteros, antes de abordar las integraciones
fraccionarias de mayor costo:

\begin{enumerate}
  \item controles analíticos lineales y verificación de simetrías;
  \item Chua PWL entero y Kalman--Fitts como controles sin PP;
  \item ampliar horizontes y brackets de MAVPD $\xi=3.1$ y PLL;
  \item MAVPD \(\xi=3.5\) y \(\xi=2.85\) para sensibilidad, retorno y borde;
  \item refinar Wu con $h/4$ y tercer solucionador, y ejecutar Muñoz \(q=0.97\);
  \item Chua PWL fraccionario y Muñoz \(q=0.996\) para comparar historia y
        conmutación;
  \item Lorenz \(q=0.995\), RF \(q=0.998\) y \texttt{c590} sólo después de pasar
        TG2 y validar el lift de memoria.
\end{enumerate}

El orden asigna primero los presupuestos bajos a controles conocidos y reserva
B3 para casos cuyos detectores, clasificadores y representaciones de memoria
ya fueron verificados.


\begin{proposalbox}
La función descriptiva, Machado, PP/FPP, las superficies críticas, las curvas
conectantes y KCC proponen semillas o regiones de interés. La integración
causal evalúa sus destinos; el seguimiento de borde estudia la separación
entre cuencas; y los métodos de conjuntos buscan certificar invariancia y
aislamiento. La campaña ha alcanzado niveles dinámicos EV--TG2/3. La
certificación de fronteras, índices y ocultedad requiere completar las
etapas posteriores.
\end{proposalbox}
